\section{Basis of Design}

\subsection{Reference standards}
The following editions form the stated basis of this worked example. They are not a
substitute for the adopted building code or amendments applicable to a real site.

\begin{longtable}{@{}p{33mm}p{115mm}@{}}
  \toprule
  \textbf{Reference} & \textbf{Application in this report} \\
  \midrule
  \endhead
  ASCE/SEI 7-22 & Minimum design loads, load combinations, and conceptual wind and
    seismic procedures. Hazard values themselves are assumed, not obtained from the
    ASCE Hazard Tool. \\
  ANSI/AISC 360-22 & LRFD resistance and stability concepts for structural steel
    members. \\
  ANSI/AISC 341-22 & Seismic detailing is identified as a required project-stage
    check; it is not completed in this example. \\
  ACI 318-19 (reapproved 2022) & Strength concepts for the preliminary reinforced
    concrete footing checks. \\
  IBC 2024 & Administrative context only; the AHJ must confirm the adopted edition
    and amendments. \\
  \bottomrule
\end{longtable}

\subsection{Design philosophy}
Strength checks use load and resistance factor design (LRFD), written generally as
\begin{equation}
  U = \sum_i \gamma_i Q_i \leq \phi R_n,
  \label{eq:lrfd}
\end{equation}
where $Q_i$ is a nominal action, $\gamma_i$ is its load factor, $R_n$ is nominal
resistance, and $\phi$ is the applicable resistance factor. Serviceability checks use
unfactored actions. Stability effects are represented by a second-order analysis
amplification stated in Section~\ref{sec:analysis}.

\subsection{Units, signs, and rounding}
Forces are reported in \si{\kilo\newton}, line loads in
\si{\kilo\newton\per\metre}, area loads in \si{\kilo\pascal}, moments in
\si{\kilo\newton\metre}, stresses in \si{\mega\pascal}, and dimensions in
\si{\milli\metre} or \si{\metre}. Downward gravity load is positive in gravity
calculations. Lateral actions may act in either sign. Displayed values are rounded;
checks use unrounded intermediate values where practical.

\subsection{Reliability and importance}
The example assumes Risk Category~II and importance factor $I_e=1.0$. No essential
facility, hazardous occupancy, special inspection, disproportionate-collapse, or
enhanced resilience requirement is assumed. These classifications are project and
jurisdiction dependent.

% =============================================================================
